% !TeX root = ../thesis.tex
% ************************** Tóm tắt *****************************
% Dùng tùy chọn `abstract' trong document class để chỉ in trang bìa và tóm tắt.
\begin{abstract}

Hỏi đáp trên ảnh tiếng Việt (Vietnamese Visual Question Answering) đòi hỏi mô hình
vừa hiểu ảnh vừa đọc được chữ tiếng Việt trong ảnh, kể cả dấu thanh. Tuy nhiên,
các vision-language model (VLM) mã nguồn mở chủ yếu được huấn luyện trên tiếng Anh,
còn việc fine-tune toàn bộ một mô hình như vậy vượt quá tài nguyên của một nhóm sinh
viên. Khóa luận này nghiên cứu cách thích nghi một VLM cho tiếng Việt chỉ bằng
parameter-efficient fine-tuning trên một GPU Google Colab.

Nhóm chọn Qwen2.5-VL-3B-Instruct làm baseline, tái lập kết quả đã công bố của mô hình
trên tập validation của TextVQA (VQA Accuracy 79,60 so với 79,3 trong báo cáo gốc),
rồi đánh giá zero-shot trên hai bộ dữ liệu tiếng
Việt: ViVQA (ảnh đời thường) và ViTextVQA (ảnh có chữ). Ở chế độ zero-shot,
mô hình chỉ đạt Exact Match (EM) 22,09 và 29,12, chủ yếu do câu trả lời sai định dạng
hoặc sai ngôn ngữ. Fine-tune bằng QLoRA 4-bit, chỉ huấn luyện 0,196\% tham số, nâng EM lên 76,64 trên ViVQA và 50,26 trên
ViTextVQA.

Cải tiến đề xuất là \textbf{QLoRA nhận biết OCR}: chữ trong ảnh được đọc bằng
PaddleOCR và VietOCR, sắp theo thứ tự đọc và đưa vào prompt cả lúc huấn luyện lẫn
lúc suy luận. Trên toàn bộ 10.028 câu test ViTextVQA, phương pháp đạt EM 53,35 và ANLS
73,89, hơn QLoRA thông thường cùng dữ liệu 3,09 điểm EM (kiểm định McNemar
$p = 6{,}7\times10^{-21}$). Mức tăng ổn định trên ba seed, và kết quả cao hơn nhiều so
với kết quả công bố tốt nhất trước đây (EM 25,48). Thí nghiệm giai thừa $2\times2$ cho thấy lợi ích chỉ xuất hiện
khi OCR có mặt cả lúc huấn luyện lẫn lúc suy luận (tương tác $+4{,}9$ EM): đóng góp
nằm ở việc dạy mô hình dùng OCR, không ở việc thêm OCR vào prompt.

Khóa luận có ablation, phân tích lỗi tự động được đối chiếu với một lượt rà soát thủ
công kiểu mù (blind), và đo hallucination bằng POPE-vi, một bộ 4.985 câu hỏi tiếng Việt
do nhóm xây dựng. Kết quả cho
thấy EM đánh giá thấp mô hình (khoảng 55\% lỗi EM là đúng ý), lỗi đọc chữ tiếng Việt
còn chiếm khoảng 17 điểm EM, và hallucination chỉ chiếm khoảng 2--4\% câu trả lời.
Mọi lần chạy dùng một GPU L4 với tối đa 8,4~GB VRAM và đều được ghi log trong một gói
tái lập. Cuối cùng, hệ thống được đóng gói thành một prototype web hiển thị chữ OCR
làm bằng chứng cho mỗi câu trả lời.

\vspace{1em}
\noindent\textbf{Từ khóa:} Vietnamese VQA, vision-language model, QLoRA,
parameter-efficient fine-tuning, OCR, scene text, hallucination.

\end{abstract}
